\documentclass[10pt,a4paper]{amsart}
\usepackage{amsmath,amssymb,amsthm,graphicx}
\usepackage{verbatim}
\newenvironment{reference}{\comment}{\endcomment}
%\newenvironment{reference}{}{}
\newenvironment{slogan}{\comment}{\endcomment}
\newenvironment{history}{\comment}{\endcomment}

% For commutative diagrams we use Xy-pic
\usepackage[all]{xy}

% We use 2cell for 2-commutative diagrams.
\xyoption{2cell}
% Theorem environments.
%
\theoremstyle{plain}
\newtheorem{theorem}[subsection]{Theorem}
\newtheorem{proposition}[subsection]{Proposition}
\newtheorem{lemma}[subsection]{Lemma}

\theoremstyle{definition}
\newtheorem{definition}[subsection]{Definition}
\newtheorem{example}[subsection]{Example}
\newtheorem{exercise}[subsection]{Exercise}
\newtheorem{situation}[subsection]{Situation}

\theoremstyle{remark}
\newtheorem{remark}[subsection]{Remark}
\newtheorem{remarks}[subsection]{Remarks}

\numberwithin{equation}{subsection}

% Macros
%
\def\lim{\mathop{\mathrm{lim}}\nolimits}
\def\colim{\mathop{\mathrm{colim}}\nolimits}
\def\Spec{\mathop{\mathrm{Spec}}}
\def\Hom{\mathop{\mathrm{Hom}}\nolimits}
\def\Ext{\mathop{\mathrm{Ext}}\nolimits}
\def\SheafHom{\mathop{\mathcal{H}\!\mathit{om}}\nolimits}
\def\SheafExt{\mathop{\mathcal{E}\!\mathit{xt}}\nolimits}
\def\Sch{\mathit{Sch}}
\def\Mor{\mathop{\mathrm{Mor}}\nolimits}
\def\Ob{\mathop{\mathrm{Ob}}\nolimits}
\def\Sh{\mathop{\mathit{Sh}}\nolimits}
\def\NL{\mathop{N\!L}\nolimits}
\def\CH{\mathop{\mathrm{CH}}\nolimits}
\def\proetale{{pro\text{-}\acute{e}tale}}
\def\etale{{\acute{e}tale}}
\def\QCoh{\mathit{QCoh}}
\def\Ker{\mathop{\mathrm{Ker}}}
\def\Im{\mathop{\mathrm{Im}}}
\def\Coker{\mathop{\mathrm{Coker}}}
\def\Coim{\mathop{\mathrm{Coim}}}

% Boxtimes
%
\DeclareMathSymbol{\boxtimes}{\mathbin}{AMSa}{"02}

%
% Macros for moduli stacks/spaces
%
\def\QCohstack{\mathcal{QC}\!\mathit{oh}}
\def\Cohstack{\mathcal{C}\!\mathit{oh}}
\def\Spacesstack{\mathcal{S}\!\mathit{paces}}
\def\Quotfunctor{\mathrm{Quot}}
\def\Hilbfunctor{\mathrm{Hilb}}
\def\Curvesstack{\mathcal{C}\!\mathit{urves}}
\def\Polarizedstack{\mathcal{P}\!\mathit{olarized}}
\def\Complexesstack{\mathcal{C}\!\mathit{omplexes}}
% \Pic is the operator that assigns to X its picard group, usage \Pic(X)
% \Picardstack_{X/B} denotes the Picard stack of X over B
% \Picardfunctor_{X/B} denotes the Picard functor of X over B
\def\Pic{\mathop{\mathrm{Pic}}\nolimits}
\def\Picardstack{\mathcal{P}\!\mathit{ic}}
\def\Picardfunctor{\mathrm{Pic}}
\def\Deformationcategory{\mathcal{D}\!\mathit{ef}}
\usepackage{cleveref}
\crefname{originallemma}{original lemma}{original lemmas}
\Crefname{originallemma}{Original lemma}{Original lemmas}
\makeatletter
\newlabel{lemma-coherent-Noetherian-quasi-coherent-sub-quotient}{{0GRB}{source}{}{sourcepointer-1}{}}
\newlabel{lemma-coherent-Noetherian-quasi-coherent-sub-quotient@cref}{{[originallemma][0][]0GRB}{[1][0][]source}{}{}{}}
\makeatother
\crefname{originalreference}{original reference}{original references}
\Crefname{originalreference}{Original reference}{Original references}
\makeatletter
\newlabel{stacks-properties-lemma-quasi-compact-stack}{{04YC}{source}{}{sourcepointer-2}{}}
\newlabel{stacks-properties-lemma-quasi-compact-stack@cref}{{[originalreference][0][]04YC}{[1][0][]source}{}{}{}}
\makeatother
\crefname{originalreference}{original reference}{original references}
\Crefname{originalreference}{Original reference}{Original references}
\makeatletter
\newlabel{algebraic-lemma-stack-presentation}{{04T5}{source}{}{sourcepointer-3}{}}
\newlabel{algebraic-lemma-stack-presentation@cref}{{[originalreference][0][]04T5}{[1][0][]source}{}{}{}}
\makeatother
\crefname{originallemma}{original lemma}{original lemmas}
\Crefname{originallemma}{Original lemma}{Original lemmas}
\makeatletter
\newlabel{lemma-coherent-presentation}{{0GRC}{source}{}{sourcepointer-4}{}}
\newlabel{lemma-coherent-presentation@cref}{{[originallemma][0][]0GRC}{[1][0][]source}{}{}{}}
\makeatother
\crefname{originalreference}{original reference}{original references}
\Crefname{originalreference}{Original reference}{Original references}
\makeatletter
\newlabel{spaces-groupoids-lemma-colimit-coherent}{{0GPS}{source}{}{sourcepointer-5}{}}
\newlabel{spaces-groupoids-lemma-colimit-coherent@cref}{{[originalreference][0][]0GPS}{[1][0][]source}{}{}{}}
\makeatother
\crefname{originalreference}{original reference}{original references}
\Crefname{originalreference}{Original reference}{Original references}
\makeatletter
\newlabel{stacks-properties-definition-type-property}{{04YG}{source}{}{sourcepointer-6}{}}
\newlabel{stacks-properties-definition-type-property@cref}{{[originalreference][0][]04YG}{[1][0][]source}{}{}{}}
\makeatother
\crefname{originalreference}{original reference}{original references}
\Crefname{originalreference}{Original reference}{Original references}
\makeatletter
\newlabel{stacks-properties-remark-list-properties-local-smooth-topology}{{04YH}{source}{}{sourcepointer-7}{}}
\newlabel{stacks-properties-remark-list-properties-local-smooth-topology@cref}{{[originalreference][0][]04YH}{[1][0][]source}{}{}{}}
\makeatother
\crefname{originalreference}{original reference}{original references}
\Crefname{originalreference}{Original reference}{Original references}
\makeatletter
\newlabel{stacks-morphisms-lemma-quasi-compact-quasi-separated-permanence}{{075S}{source}{}{sourcepointer-8}{}}
\newlabel{stacks-morphisms-lemma-quasi-compact-quasi-separated-permanence@cref}{{[originalreference][0][]075S}{[1][0][]source}{}{}{}}
\makeatother
\crefname{originalreference}{original reference}{original references}
\Crefname{originalreference}{Original reference}{Original references}
\makeatletter
\newlabel{spaces-morphisms-lemma-finite-presentation-noetherian}{{04ZL}{source}{}{sourcepointer-9}{}}
\newlabel{spaces-morphisms-lemma-finite-presentation-noetherian@cref}{{[originalreference][0][]04ZL}{[1][0][]source}{}{}{}}
\makeatother
\title{On a Noetherian algebraic stack every qc module is a filtered colimit of coherent submodules}
\author{The Stacks Project Authors}
\date{stacks-project:stacks-cohomology.tex:local-snapshot-20261012; GNU FDL 1.2 or later}
\begin{document}
\maketitle
\noindent\textit{Editorial note.} All delivered mathematical text is the author's own, copied verbatim from the pinned Stacks Project chapter file stacks-cohomology.tex; no statement or proof step was rewritten, reordered or removed.\par\medskip
\noindent\textit{Editorial note.} The counted span cites the lemmas recorded in assembly.external\_references with canonical tags as navigation pointers.\par\medskip
\noindent\textit{Allowed standard prerequisites (editorial).}\begin{itemize}
\item Coherent submodules of qc modules. Cohomology of Algebraic Stacks, Lemma stacks-cohomology-lemma-coherent-Noetherian-quasi-coherent-sub-quotient. Images and sums of coherent submodules are coherent.
\item Groupoid presentation of QCoh and Coh. Cohomology of Algebraic Stacks, Lemma stacks-cohomology-lemma-coherent-presentation; Algebraic Stacks, Lemma algebraic-lemma-stack-presentation. Identifies the categories with groupoid modules.
\item Filtered colimits of coherent modules on groupoids. Groupoids in Spaces, Lemma spaces-groupoids-lemma-colimit-coherent. The groupoid-level statement being transported.
\end{itemize}\medskip
\ifdefined\hypertarget\hypertarget{sourcepointer-1}{}\fi
\noindent\textit{Source reference.} Original lemma 0GRB: sources/stacks-chow/tags.txt, source lines 20891--20891. Cross-chapter standard prerequisite cited by the author inside the counted span; it is used as a pointer and its own proof is not part of this item.\par\medskip
\ifdefined\hypertarget\hypertarget{sourcepointer-2}{}\fi
\noindent\textit{Source reference.} Original reference 04YC: sources/stacks-chow/tags.txt, source lines 6425--6425. Cross-chapter standard prerequisite cited by the author inside the counted span; it is used as a pointer and its own proof is not part of this item.\par\medskip
\ifdefined\hypertarget\hypertarget{sourcepointer-3}{}\fi
\noindent\textit{Source reference.} Original reference 04T5: sources/stacks-chow/tags.txt, source lines 6243--6243. Cross-chapter standard prerequisite cited by the author inside the counted span; it is used as a pointer and its own proof is not part of this item.\par\medskip
\ifdefined\hypertarget\hypertarget{sourcepointer-4}{}\fi
\noindent\textit{Source reference.} Original lemma 0GRC: sources/stacks-chow/tags.txt, source lines 20892--20892. Cross-chapter standard prerequisite cited by the author inside the counted span; it is used as a pointer and its own proof is not part of this item.\par\medskip
\ifdefined\hypertarget\hypertarget{sourcepointer-5}{}\fi
\noindent\textit{Source reference.} Original reference 0GPS: sources/stacks-chow/tags.txt, source lines 20837--20837. Cross-chapter standard prerequisite cited by the author inside the counted span; it is used as a pointer and its own proof is not part of this item.\par\medskip
\ifdefined\hypertarget\hypertarget{sourcepointer-6}{}\fi
\noindent\textit{Source reference.} Original reference 04YG: sources/stacks-chow/tags.txt, source lines 6429--6429. Cross-chapter standard prerequisite cited by the author inside the counted span; it is used as a pointer and its own proof is not part of this item.\par\medskip
\ifdefined\hypertarget\hypertarget{sourcepointer-7}{}\fi
\noindent\textit{Source reference.} Original reference 04YH: sources/stacks-chow/tags.txt, source lines 6430--6430. Cross-chapter standard prerequisite cited by the author inside the counted span; it is used as a pointer and its own proof is not part of this item.\par\medskip
\ifdefined\hypertarget\hypertarget{sourcepointer-8}{}\fi
\noindent\textit{Source reference.} Original reference 075S: sources/stacks-chow/tags.txt, source lines 9141--9141. Cross-chapter standard prerequisite cited by the author inside the counted span; it is used as a pointer and its own proof is not part of this item.\par\medskip
\ifdefined\hypertarget\hypertarget{sourcepointer-9}{}\fi
\noindent\textit{Source reference.} Original reference 04ZL: sources/stacks-chow/tags.txt, source lines 6469--6469. Cross-chapter standard prerequisite cited by the author inside the counted span; it is used as a pointer and its own proof is not part of this item.\par\medskip
\begin{lemma}
\label{lemma-directed-colimit-coherent}
Let $\mathcal{X}$ be a Noetherian algebraic stack. Every quasi-coherent
$\mathcal{O}_\mathcal{X}$-module is the filtered colimit of its coherent
submodules.
\end{lemma}

\begin{proof}
Let $\mathcal{F}$ be a quasi-coherent $\mathcal{O}_\mathcal{X}$-module.
If $\mathcal{G}, \mathcal{H} \subset \mathcal{F}$ are coherent
$\mathcal{O}_\mathcal{X}$-submodules then the image of
$\mathcal{G} \oplus \mathcal{H} \to \mathcal{F}$ is another
coherent $\mathcal{O}_\mathcal{X}$-submodule which contains both of them,
see Lemma \ref{lemma-coherent-Noetherian-quasi-coherent-sub-quotient}.
In this way we see that the system is directed.
Hence it now suffices to show that $\mathcal{F}$ can be written as
a filtered colimit of coherent modules, as then we can take the
images of these modules in $\mathcal{F}$ to conclude there are
enough of them.

\medskip\noindent
Let $U$ be an affine scheme and $U \to \mathcal{X}$
a surjective smooth morphism (Properties of Stacks, Lemma
\ref{stacks-properties-lemma-quasi-compact-stack}).
Set $R = U \times_\mathcal{X} U$ so that $\mathcal{X} = [U/R]$
as in Algebraic Stacks, Lemma \ref{algebraic-lemma-stack-presentation}.
By Lemma \ref{lemma-coherent-presentation} we have
$\QCoh(\mathcal{O}_X) = \QCoh(U, R, s, t, c)$ and
$\textit{Coh}(\mathcal{O}_X) = \textit{Coh}(U, R, s, t, c)$.
In this way we reduce to the problem of proving the corresponding thing for
$\QCoh(U, R, s, t, c)$. This is
Groupoids in Spaces, Lemma \ref{spaces-groupoids-lemma-colimit-coherent};
we check its assumptions in the next paragraph.

\medskip\noindent
We urge the reader to skip the rest of the proof.
The affine scheme $U$ is Noetherian; this follows from our definition
of $\mathcal{X}$ being locally Noetherian, see
Properties of Stacks, Definition
\ref{stacks-properties-definition-type-property} and
Remark \ref{stacks-properties-remark-list-properties-local-smooth-topology}.
The projection morphisms $s, t : R \to U$ are smooth (see reference
given above)
and quasi-separated and quasi-compact (Morphisms of Stacks, Lemma
\ref{stacks-morphisms-lemma-quasi-compact-quasi-separated-permanence}).
In particular, $R$ is a quasi-compact and quasi-separated algebraic space
smooth over $U$ and hence Noetherian
(Morphisms of Spaces, Lemma
\ref{spaces-morphisms-lemma-finite-presentation-noetherian}).
\end{proof}

\end{document}
